\CWHeader{Лабораторная работа \textnumero 1}

\CWProblem{
Необходимо разработать программную библиотеку, реализующую сбалансированное дерево поиска, и на её основе создать программу-словарь, обрабатывающую команды из стандартного ввода до его окончания.

{\bfseries Вариант структуры данных:} Красно-чёрное дерево.

{\bfseries Ограничения:} время работы --- до 15 секунд; оперативная память --- до 512~МБ.

{\bfseries Формат данных:} ключом является регистронезависимая последовательность букв английского алфавита длиной до 256 символов; значением --- беззнаковое 64-битное целое число от $0$ до $2^{64} - 1$.

{\bfseries Поддерживаемые команды:}
\begin{itemize}
    \item \texttt{+ word value} --- добавить пару в словарь (\texttt{OK} или \texttt{Exist});
    \item \texttt{- word} --- удалить слово из словаря (\texttt{OK} или \texttt{NoSuchWord});
    \item \texttt{word} --- поиск слова (\texttt{OK: value} или \texttt{NoSuchWord});
    \item \texttt{! Save /path/to/file} --- сохранить словарь в компактном бинарном формате;
    \item \texttt{! Load /path/to/file} --- загрузить словарь из файла с проверкой формата и сохранением текущего словаря при ошибке.
\end{itemize}

При системных ошибках или ошибках формата программа выводит строку с префиксом \texttt{ERROR:} и описанием на английском языке.

{\bfseries Исследовательская часть:} провести анализ скорости выполнения и потребления оперативной памяти с использованием утилит \texttt{gprof}, \texttt{Valgrind} (\texttt{Memcheck}, \texttt{Massif}) и \texttt{/usr/bin/time -v}; выявить и устранить недочёты реализации; сравнить метрики исходной и исправленной версий программы.
}
\pagebreak
